\subsubsection{正交切片与判定闭合}\label{subsubsec:typed-address-biaxial-completion-orthogonal-trislice-decision-closure}
Jensen 半径族、Toeplitz--PSD 端点热核与共动 Fourier--Laplace 层析分别给出三类读出：盘内零点、端点载荷与边界层缺陷测度。判定层必须进一步说明这些读出是否来自同一个解析对象，哪些读出能够粘接，哪些预算不足只能给出 \(\mathrm{NULL}\)，以及何时证书矛盾应输出 \(\mathrm{FAIL}\)。相应的最小闭包由三条正交切片组成：
\[
\text{visible modulus slice},
\qquad
\text{Blaschke zero slice},
\qquad
\text{endpoint singular/load slice}.
\]
它们分别记录可见模长、盘内零点与端点奇异载荷，且互不抵偿。

\paragraph{解析对象的三因子分解。}
令 \(F\) 是单位圆盘 \(\mathbb D\) 上的非零解析对象，处于 Hardy/Smirnov 型口径中。设边界极限满足
\[
h(\theta):=\log|F^*(e^{i\theta})|\in L^1(\partial\mathbb D).
\]
盘内零点记为
\[
\mathcal Z:=\sum_j m_j\delta_{a_j},
\qquad
a_j\in\mathbb D,
\]
并满足 Blaschke 条件
\[
\sum_jm_j(1-|a_j|)<\infty.
\]
端点或边界奇异载荷记为有限正奇异测度 \(\mu_\partial\)。内外分解给出
\[
\boxed{
F(w)=
e^{i\alpha}
O_h(w)
B_{\mathcal Z}(w)
S_{\mu_\partial}(w).
}
\]
其中
\[
O_h(w)
=
\exp\left\{
\frac1{2\pi}
\int_0^{2\pi}
\frac{e^{i\theta}+w}{e^{i\theta}-w}
h(\theta)\,d\theta
\right\},
\]
\[
B_{\mathcal Z}(w)
=
\prod_j
\left(
\frac{|a_j|}{a_j}
\frac{a_j-w}{1-\overline a_j w}
\right)^{m_j},
\]
零点 \(a_j=0\) 时取相应规范形式，并且
\[
S_{\mu_\partial}(w)
=
\exp\left\{
-\int_{\partial\mathbb D}
\frac{\zeta+w}{\zeta-w}\,d\mu_\partial(\zeta)
\right\}.
\]
定义
\[
\boxed{
\mathsf{TriSlice}(F):=(h,\mathcal Z,\mu_\partial).
}
\]
全局相位 \(e^{i\alpha}\) 是 \(U(1)\) 规范自由度，单独记入 \(\mathsf{PhaseGauge}\)。

\begin{theorem}[正交三切片唯一性]\label{thm:typed-address-biaxial-completion-orthogonal-trislice-uniqueness}
在上述 admissible 类中，若非零解析对象 \(F,G\) 满足
\[
h_F=h_G,
\qquad
\mathcal Z_F=\mathcal Z_G,
\qquad
\mu_{\partial,F}=\mu_{\partial,G},
\]
则存在 \(\alpha\in\mathbb R\)，使得
\[
F=e^{i\alpha}G.
\]
\end{theorem}
\begin{proof}
由内外分解，
\[
F=e^{i\alpha_F}O_hB_{\mathcal Z}S_{\mu_\partial},
\qquad
G=e^{i\alpha_G}O_hB_{\mathcal Z}S_{\mu_\partial}.
\]
相除得到 \(F/G=e^{i(\alpha_F-\alpha_G)}\)。证毕。
\end{proof}

定义投影
\[
\Pi_{\rm mod}(F):=h=\log|F^*|,
\qquad
\Pi_{\rm zero}(F):=\mathcal Z,
\qquad
\Pi_{\rm end}(F):=\mu_\partial.
\]
边界模长由 Hardy 投影、有理核与有限 Fourier 模态恢复；盘内零点由 Jensen 缺陷半径族恢复：
\[
J_F(r)
=
\frac1{2\pi}
\int_0^{2\pi}\log|F(re^{i\theta})|\,d\theta
-\log|F(0)|
=
\sum_{|a_j|<r}m_j\log\frac r{|a_j|}.
\]
在无零半径处，
\[
\frac{d}{d\log r}J_F(r)=N_F(r).
\]
端点载荷由 Toeplitz--PSD、Carath\'eodory 正性、端点热核或近端 Poisson 核恢复。若 \(K_N(\eta,\zeta)\) 是以 \(\eta\in\partial\mathbb D\) 为中心的正近似单位核，则
\[
H_N(\eta):=\int_{\partial\mathbb D}K_N(\eta,\zeta)\,d\mu_\partial(\zeta)
\longrightarrow
\mu_\partial(\{\eta\})
\]
在标准局部化条件下成立。

\paragraph{三切片不可替代。}
Blaschke 因子与奇异内因子满足
\[
|B_{\mathcal Z}^*(e^{i\theta})|=1,
\qquad
|S_{\mu_\partial}^*(e^{i\theta})|=1
\]
对 Lebesgue 几乎处处边界点成立，因此 \(\Pi_{\rm mod}\) 对内因子盲。外函数的圆周平均满足调和平均恒等式；奇异内因子满足
\[
\frac1{2\pi}\int_0^{2\pi}
\log|S_{\mu_\partial}(re^{i\theta})|\,d\theta
=
\log|S_{\mu_\partial}(0)|.
\]
故 Jensen 缺陷的半径导数只读盘内零点。端点载荷则可在边界模长几乎处处不变、盘内零点不变的情况下改变：若
\[
\mu_\partial'=\mu_\partial+\tau\delta_\eta,
\qquad \tau>0,
\]
则零点切片不变，边界模长仍为 \(1\) 的内因子模长，而端点热核在 \(\eta\) 处的极限增加 \(\tau\)。

\begin{theorem}[三预算不可抵偿]\label{thm:typed-address-biaxial-completion-trislice-budgets-noncompensating}
存在三组对象，使得任意两个切片相同而第三个切片不同。因此，模长、零点与端点载荷三类预算互不抵偿。
\end{theorem}
\begin{proof}
固定 \(h,\mu_\partial\)，取两个不同 Blaschke 因子，得到零点切片不同而其余两切片相同。固定 \(h,\mathcal Z\)，把 \(\mu_\partial\) 改为 \(\mu_\partial+\tau\delta_\eta\)，得到端点切片不同而其余两切片相同。固定 \(\mathcal Z,\mu_\partial\)，取两个不同的 \(h\in L^1\)，得到模长切片不同而其余两切片相同。证毕。
\end{proof}

\begin{definition}[可见模长证书]\label{def:typed-address-biaxial-completion-modulus-slice-cert}
称 \(\mathsf{ModulusSliceCert}_M\) 为可见模长证书，若它包含边界分割 \(\{I_q\}_{q=1}^Q\)、区间包络 \(\{[h_q^-,h_q^+]\}_{q=1}^Q\)、boundary sampling ledger、Fourier tail ledger、outer reconstruction ledger 与 round hash，并证明
\[
\|h-h_M\|_{L^1}\le\varepsilon_{\rm mod}.
\]
若给定范数 \(\mathcal H\) 下外函数重构满足
\[
\|O_h-O_{h_M}\|_{\mathcal H}
\le
C_{\rm out}\varepsilon_{\rm mod},
\]
则该误差进入外因子预算。若边界分辨率不足以稳定确定 \(h\)，输出 \(\mathsf{ModulusResolutionNull}\)。
\end{definition}

\begin{definition}[零点切片证书]\label{def:typed-address-biaxial-completion-zero-slice-cert}
称 \(\mathsf{ZeroSliceCert}\) 为零点切片证书，若它包含半径 \(\{\rho_j\}_{j=1}^M\)、区间 \(\{[J(\rho_j)^-,J(\rho_j)^+]\}_{j=1}^M\)、Jensen circle average ledger、log-radius derivative ledger、零点盒 \(\{([a_\ell],m_\ell)\}_{\ell=1}^L\)、zero-box separation ledger、residual ledger 与 round hash。验证器核验
\[
J(\rho_j)
=
\frac1{2\pi}
\int_0^{2\pi}\log|F(\rho_je^{i\theta})|\,d\theta
-\log|F(0)|
\]
以及差分导数与提交计数区间相容。若零点盒贴近半径阈值，输出 \(\mathsf{RadiusThresholdNull}\)；若当前前缀无法分离零点盒，输出 \(\mathsf{ZeroCollisionNull}\)。
\end{definition}

\begin{definition}[端点载荷证书]\label{def:typed-address-biaxial-completion-endpoint-load-cert}
称 \(\mathsf{EndpointLoadCert}_{N,\eta}\) 为端点载荷证书，若它包含矩区间 \(\{[m_k^-,m_k^+]\}_{|k|\le N}\)、Toeplitz PSD interval cert、heat kernel ledger、区间 \([H_N(\eta)^-,H_N(\eta)^+]\)、endpoint tail ledger 与 round hash，其中
\[
m_k=\int_{\partial\mathbb D}\zeta^{-k}\,d\mu_\partial(\zeta)
\]
且
\[
T_N=(m_{j-k})_{0\le j,k\le N}\succeq0.
\]
若核族给出近似单位，则
\[
H_N(\eta)=\mu_\partial(\{\eta\})+O(\varepsilon_N).
\]
谱隙不足或热核阶数不够时输出 \(\mathsf{EndpointHeatDepthNull}\)。Toeplitz 区间无法证明 PSD 时输出 \(\mathsf{ToeplitzPSDNull}\)，若存在负谱见证则输出 \(\mathsf{ToeplitzPSDFail}\)。
\end{definition}

\begin{definition}[正交三切片证书]\label{def:typed-address-biaxial-completion-orthogonal-trislice-cert}
称
\[
\mathsf{OrthogonalTriSliceCert}
\]
为三切片总证书，若它由
\[
\mathsf{ModulusSliceCert},\quad
\mathsf{ZeroSliceCert},\quad
\mathsf{EndpointLoadCert},\quad
\mathsf{PhaseGaugeLedger},\quad
\mathsf{FactorizationResidualLedger},\quad
\mathsf{OrthogonalBudgetLedger},\quad
\mathsf{TranscriptHash}
\]
组成，并证明重构对象
\[
F_{\rm rec}(w)
=
e^{i\alpha}O_{\widehat h}(w)
B_{\widehat{\mathcal Z}}(w)
S_{\widehat\mu}(w)
\]
满足
\[
\|F-F_{\rm rec}\|_{\mathcal H}\le\varepsilon_{\rm tri},
\]
或在只读协议下满足所有投影一致性约束。总误差账本为
\[
\varepsilon_{\rm tri}
=
\varepsilon_{\rm mod}
+\varepsilon_{\rm zero}
+\varepsilon_{\rm end}
+\varepsilon_{\rm phase}
+\varepsilon_{\rm fact}.
\]
该式是误差求和，不表示预算可互相替代。
\end{definition}

\paragraph{三值读出。}
令地址对象为
\[
a=(L,\theta,b,\mathsf{Mode})
\]
并令预算包为
\[
B=(\Delta_\rho,b,N,E_\tau,\Omega,\Lambda,L_{\rm Hankel},C_{\rm cond},\varepsilon).
\]
定义统一读出偏函数
\[
\boxed{
\mathsf{ReadUS}_B(a)
\in
\{\mathsf{VAL}(v),\mathsf{NULL}(\tau,W),\mathsf{FAIL}(R)\}.
}
\]
\(\mathsf{VAL}(v)\) 表示值被稳定读出；\(\mathsf{NULL}(\tau,W)\) 表示协议内不可定义或不可决定，并给出类型与见证；\(\mathsf{FAIL}(R)\) 表示提交证书与验证规则矛盾。因此 \(\mathsf{NULL}\ne\mathsf{FAIL}\)。

NULL 见证分为六类：
\[
\mathsf{DomainOutNull},\quad
\mathsf{LocalEmptyNull},\quad
\mathsf{CechGluingNull},\quad
\mathsf{HardThresholdNull},\quad
\mathsf{CollisionNull},\quad
\mathsf{ValueMarginNull}.
\]
其中 Čech 粘接障碍由闭路 holonomy 给出。若局部截面 \(\sigma_i\in\mathcal F(U_i)\) 在重叠上满足
\[
\sigma_i=g_{ij}\cdot\sigma_j,
\]
严格粘接要求存在局部规范 \(h_i\) 使
\[
g_{ij}=h_i h_j^{-1}.
\]
若闭路
\[
\Gamma=(i_0,i_1,\ldots,i_m=i_0)
\]
满足
\[
\operatorname{Hol}_\Gamma
:=
g_{i_0i_1}g_{i_1i_2}\cdots g_{i_{m-1}i_m}
\ne e,
\]
则输出 \(\mathsf{CechGluingNull}\)。

\begin{definition}[非 NULL 证书]\label{def:typed-address-biaxial-completion-not-null-cert}
称 \(\mathsf{NotNullCert}_B(a)\) 为非 NULL 证书，若它包含 domain membership cert、local section cert、Čech trivialization cert、OrthogonalTriSliceCert、budget pass cert、collision separation cert、value margin cert 与 round hash。它证明 \(a\in\mathsf{AdmAddr}_B\)、局部截面存在、\(g_{ij}=h_i h_j^{-1}\)、三切片均通过、预算达标、候选 lift 已分离，并给出值区间
\[
v\in[v^-,v^+],
\qquad v^+-v^-<\varepsilon.
\]
阈值谓词 \(v\le\Theta\) 还必须有严格 margin：
\[
v^+<\Theta
\quad\text{or}\quad
v^->\Theta.
\]
若缺少 margin，则输出 \(\mathsf{ValueMarginNull}\)。
\end{definition}

\begin{theorem}[NULL 与非 NULL 的有限判定闭合]\label{thm:typed-address-biaxial-completion-null-not-null-finite-closure}
固定有限协议包 \(B\)、有限覆盖 \(\mathcal U\)、有限证书语言和有限误差阈值。若验证器包含 domain、local section、Čech、tri-slice、budget、collision 与 margin 七个 gate，则对任意地址 \(a\)，\(\mathsf{ReadUS}_B(a)=\mathsf{NULL}\) 或 \(\mathsf{ReadUS}_B(a)=\mathsf{VAL}(v)\) 均有有限证书；若提交证书内部矛盾，则输出 \(\mathsf{FAIL}\)。
\end{theorem}
\begin{proof}
固定 \(B\) 后，每个 gate 均是有限验证对象。域检查失败给出 \(\mathsf{DomainOutNull}\)；局部片不可截面化给出 \(\mathsf{LocalEmptyNull}\)；局部截面存在但 Čech residual 具有非平凡闭路 holonomy，给出 \(\mathsf{CechGluingNull}\)；三切片或预算不足给出 \(\mathsf{HardThresholdNull}\)；存在同可见 lift 而目标值不同，给出 \(\mathsf{CollisionNull}\)；阈值无裕度给出 \(\mathsf{ValueMarginNull}\)。若这些障碍均不存在，则域内、局部截面存在、粘接平凡、三切片通过、预算达标、碰撞分离且 margin 成立，故值由有限证书唯一确定。证书声称通过而复算不支持时，输出 \(\mathsf{FAIL}\)。证毕。
\end{proof}

\begin{definition}[正交判定融合证书]\label{def:typed-address-biaxial-completion-orthogonal-decision-fusion-cert}
最终判定证书记为
\[
\mathsf{OrthogonalDecisionFusionCert}.
\]
其数据为
\[
\left(
a,B,\mathcal P,
\mathsf{OrthogonalTriSliceCert},
\mathsf{CechTrivializationCert}\ \text{or}\ \mathsf{CechObstacleCert},
\mathsf{NotNullCert}\ \text{or}\ \mathsf{NullCert},
\mathsf{DecisionMarginLedger},
\mathsf{RejectLedger},
\mathsf{TranscriptHash}
\right).
\]
其中 \(\mathsf{NullCert}\) 必须属于上列六类 NULL 之一。
\end{definition}

验证器 \(\mathsf{VerifyOrthogonalDecisionFusion}\) 检查地址类型、局部截面、Čech 粘接、三切片、三端预算、碰撞分离以及目标谓词区间。若目标谓词在三端误差盒
\[
\mathcal H_{\rm tri}
=
\mathcal H_{\rm mod}\times\mathcal H_{\rm zero}\times\mathcal H_{\rm end}
\]
中不跨阈值，则输出相应判定；若跨阈值，则输出 \(\mathsf{ValueMarginNull}\)。

\begin{theorem}[正交判定融合证书的可靠性]\label{thm:typed-address-biaxial-completion-orthogonal-decision-fusion-soundness}
若
\[
\mathsf{VerifyOrthogonalDecisionFusion}
(\mathsf{OrthogonalDecisionFusionCert})
=
\mathrm{PASS},
\]
则目标谓词 \(\mathcal P(F)\) 在当前 admissible 类与预算 \(B\) 下成立。若输出 \(\mathsf{NULL}(\tau,W)\)，则 \(W\) 是当前协议内无法给出唯一 PASS/FAIL 的有限见证；若输出 \(\mathsf{FAIL}(R)\)，则提交证书与验证规则冲突。
\end{theorem}
\begin{proof}
PASS 表示地址在 typed domain 内，局部截面存在，Čech residual 已平凡化，三切片证书分别通过，预算达标且碰撞已分离，并且目标谓词在整个三端误差盒中具有同一真值。故 \(\mathcal P(F)\) 成立。NULL 情形由某个 gate 的有限见证触发，表示当前协议不能唯一读出目标值；FAIL 情形才表示等式、PSD、holonomy、预算或残差不被验证器接受。证毕。
\end{proof}

前述端点热核融合证书嵌入 \(\mathsf{EndpointLoadCert}\)；共动层析融合证书为 \(\mathsf{ZeroSliceCert}\) 与 \(\mathsf{EndpointLoadCert}\) 提供局部增强。最终判定仍需要 Čech trivialization 或 Čech obstacle：局部数据即使全部正确，也未必能粘接为全局对象。

若过渡群为 \(U(1)\)，则 residual 写为
\[
g_{ij}=e^{i\varphi_{ij}},
\]
闭路 holonomy 为
\[
\operatorname{Hol}_\Gamma
=
\exp\left(i\sum_{(i,j)\in\Gamma}\varphi_{ij}\right).
\]
若 \(\operatorname{Hol}_\Gamma\ne1\)，则不能同时选取局部相位消去所有 residual。该情形输出 \(\mathsf{CechGluingNull}\)，或在允许射影化的协议中输出 \(\mathsf{ProjectiveLiftRequired}\)。这把相位接口解释为粘接 residual 的可见投影。
